% Generated draft: requires mathematical and source review.
\documentclass[11pt]{article}
\usepackage[margin=25mm]{geometry}
\usepackage{fontspec}
\setmainfont{Latin Modern Roman}
\usepackage{amsmath,amssymb,mathtools,unicode-math}
\setmathfont{Latin Modern Math}
\usepackage{enumitem,needspace,xurl,hyperref}
\setlength{\parindent}{0pt}
\setlength{\parskip}{5pt}
\setlength{\emergencystretch}{3em}
\newcommand{\N}{\mathbb{N}}
\newcommand{\Z}{\mathbb{Z}}
\newcommand{\Q}{\mathbb{Q}}
\newcommand{\R}{\mathbb{R}}
\newcommand{\C}{\mathbb{C}}
\newcommand{\F}{\mathbb{F}}
\newcommand{\A}{\mathbb{A}}
\newcommand{\PP}{\mathbb{P}}
\newcommand{\E}{\mathbb{E}}
\newcommand{\eps}{\varepsilon}
\newcommand{\dd}{\,\mathrm{d}}
\begin{document}
\section*{UnsolvedMath ID: 30000694}
\textbf{Source problem code:} catalogue entry
\subsection{Definitions and mathematical statement}
Let \(d\geq1\), let \(c_1,\ldots,c_m\in\R^d\) with \(1\leq m<\infty\), and set
\[
P=\bigcap_{i=1}^m\overline B(c_i,1)
=\{x\in\R^d:\|x-c_i\|_2\leq1\text{ for all }i\}.
\]
The source asks for a facial theory of these intersections of congruent Euclidean balls in arbitrary dimension. It does not impose nonempty interior, irredundancy of the balls, distinct centers, or general position; such hypotheses must be specified before using a standard ball-polyhedron theory.

The intended faces are presumably spherical boundary patches and their incidences, not just ordinary convex faces defined by supporting hyperplanes. For a full-dimensional intersection, strict convexity makes the latter unsuitable as a finite polyhedral-style face structure. One possible source of strata is the active-constraint data
\[
F_I=P\cap\bigcap_{i\in I}\partial B(c_i,1),\qquad
S_I=\{x\in P:\|x-c_i\|=1\ (i\in I),\ \|x-c_j\|<1\ (j\notin I)\}.
\]
A facial definition would need to say whether faces are these sets, closures of their connected components, maximal spherical patches, or another structure. The draft does not assert that all such sets are cells or that their incidence poset is a lattice. Redundant or repeated defining balls also need to be treated consistently if the structure is to depend on \(P\) rather than its chosen presentation.

The research task is to determine the allowable face types, dimensions, topology, and incidence behavior under the intended definition, and to identify the hypotheses under which a polyhedral-style complex or combinatorial characterization exists. No specific classification, incidence axiom, or face-number inequality is stated in the extract, so this remains a structural programme rather than a uniquely quantified conjecture.

The supplied assessment checked 2026-08-21 reports substantial results for three-dimensional standard ball-polyhedra and other structural subclasses, but an incomplete arbitrary-dimensional theory. Standardness or nondegeneracy from those subclasses is not imported into the general question.
\subsection{Short English statement}
Describe the spherical-patch facial structure of finite intersections of unit balls in arbitrary dimension. The source must fix what counts as a face and how degenerate or redundant ball presentations are handled; ordinary convex faces are a different notion.
\subsection{Sources}
\begin{itemize}
\item Local UnsolvedMath registry, permanent ID 30000694, catalogue entry, \emph{Facial Structure of Intersections of Unit Balls}.
\item \emph{Geometric and Topological Combinatorics} (2007), Oberwolfach Reports, pp. 195--272, DOI \url{https://doi.org/10.4171/owr/2007/04}.
\item Supplied reference \emph{Ball-Polyhedra}: \url{https://arxiv.org/abs/1110.4329}.
\item Supplied reference ``selected topics on intersections of balls'': \url{https://arxiv.org/abs/2411.10302}.
\item Supplied literature triage checked 2026-08-21; exact standardness and special-case hypotheses are not reproduced.
\end{itemize}
\end{document}
